% cluster_definition.png: how Ansys finds clusters, drawn on a triangular lattice. The sites come from analysis_schematics.py.
% Make it in this folder with, for any scratch folder T:
%   python3 analysis_schematics.py T
%   TEXINPUTS=.:T: pdflatex -output-directory=T fig_cluster_definition.tex
%   python3 pdf_to_png.py T/fig_cluster_definition.pdf ../cluster_definition.png
\documentclass[11pt,tikz,border=1.5pt]{standalone}
% Shared preamble of all figures. Every fig_*.tex is compiled on its own, see the comment at its top; the fonts are those of the matplotlib
% style xzc_sans (amsmath, amssymb, amsfonts, sfmath), so that line drawings and data plots look the same.
\usepackage{amsmath,amssymb,amsfonts}
\usepackage{sfmath}
\renewcommand{\familydefault}{\sfdefault}
\usepackage{graphicx}
\usetikzlibrary{arrows.meta,calc,positioning,decorations.pathreplacing,backgrounds,shapes.geometric}
% species: Al gray, Mg orange, Zn blue; red marks the vacancy and the jumping atom
\definecolor{cAl}{HTML}{B8BCC2}
\definecolor{cMg}{HTML}{E08A2E}
\definecolor{cZn}{HTML}{3F78B5}
\definecolor{cOne}{HTML}{C8423B}
\definecolor{cTwo}{HTML}{3F78B5}
\definecolor{cThree}{HTML}{3E9B5C}
\definecolor{cFour}{HTML}{E0A21B}
\definecolor{cGray}{HTML}{7E858F}
\definecolor{cDark}{HTML}{2B2B2B}
\newcommand{\panel}[1]{\textbf{(#1)}}
\newcommand{\DM}{D_{\mathrm{MEP}}}
\newcommand{\KI}{K_{\mathrm{I}}}
\newcommand{\KF}{K_{\mathrm{F}}}
\newcommand{\dE}{\Delta E}
\newcommand{\Ea}{\Delta E_{\mathrm{a}}}
\newcommand{\ea}{\Delta e_{\mathrm{a}}}
\newcommand{\Ebind}{E_{\mathrm{bind}}^{\mathrm{site}}}
% a ball with a soft highlight, like the rendered spheres
\makeatletter
\pgfdeclareradialshading[tikz@ball]{softball}{\pgfpoint{-9bp}{9bp}}{%
  color(0bp)=(tikz@ball!22!white); color(6bp)=(tikz@ball!62!white); color(13bp)=(tikz@ball!96!white);
  color(20bp)=(tikz@ball!92!black); color(27bp)=(tikz@ball!74!black); color(50bp)=(tikz@ball!55!black)}
\makeatother
\tikzset{
  every picture/.append style={line cap=round, line join=round},
  every node/.append style={font=\small},
  img/.style={inner sep=0pt, anchor=south west},
  lab/.style={font=\normalsize, inner sep=1pt, anchor=north west},
  above lab/.style={font=\normalsize, inner sep=1pt, anchor=south west},
  capt/.style={align=center, anchor=north},
  atom/.style={circle, ball color=#1, shading=softball, minimum size=12pt, inner sep=0pt},
  atom/.default=cAl,
  vac/.style={circle, draw=cOne, line width=1.1pt, fill=white, minimum size=11.5pt, inner sep=0pt},
  arr/.style={-{Stealth[length=5pt]}, line width=0.8pt},
  dbl/.style={{Stealth[length=4pt]}-{Stealth[length=4pt]}, line width=0.8pt},
  jump/.style={arr, cOne, line width=1.2pt},
  block/.style={-{Triangle[length=8pt, width=13pt]}, line width=5pt, cTwo!38},
  cell axis/.style={-{Triangle[length=6pt, width=5.5pt]}, line width=1.2pt, black},
  dir/.style={fill=white, fill opacity=0.75, text opacity=1, font=\footnotesize, inner sep=1.5pt},
}
\graphicspath{{img/}}
% data plots are included at one scale: \plot{name}{x}{y}{file}
\newcommand{\plotscale}{0.6}
\newcommand{\plot}[4]{\node[img] (#1) at (#2,#3) {\includegraphics[scale=\plotscale]{#4}};}
\newcommand{\plotlabel}[2]{\node[above lab] at ($(#1.north west)+(-0.05,0.02)$) {\panel{#2}};}
% legend of the species, first entry at (#1,#2): Al, Mg, Zn (and the vacancy)
\tikzset{leg/.style={anchor=base west, inner sep=2pt, yshift=-3.2pt}}
\newcommand{\specieslegend}[2]{%
  \foreach \lx/\lc/\lt in {0/cAl/Al, 1.2/cMg/Mg, 2.4/cZn/Zn} {
    \node[atom=\lc, minimum size=10pt] at ({#1+\lx},#2) {}; \node[leg] at ({#1+\lx+0.2},#2) {\lt};}}
\newcommand{\vacancylegend}[2]{\node[vac, minimum size=9.5pt] at (#1,#2) {}; \node[leg] at ({#1+0.2},#2) {vacancy};}

\input{cluster_definition_sites}
\begin{document}
\begin{tikzpicture}[x=0.56cm, y=0.56cm,
    al/.style={atom=cAl, minimum size=11.4pt}, mg/.style={atom=cMg, minimum size=13.2pt}, zn/.style={atom=cZn, minimum size=11.4pt},
    va/.style={vac, minimum size=11.2pt},
    joined/.style={circle, draw=cThree, line width=1.8pt, minimum size=12pt, inner sep=0pt},
    bond/.style={cDark, line width=2.2pt},
    rim/.style={cTwo!75, line width=16pt}, region/.style={cTwo!20, line width=14.6pt},
    number/.style={font=\footnotesize, inner sep=0pt},
    pointer/.style={black!70, line width=0.6pt}]
  \newcommand{\lines}[2]{\foreach \x/\y/\u/\v in #1 \draw[#2] (\x,\y) -- (\u,\v);}
  \newcommand{\nodes}[2]{\foreach \x/\y in #1 \node[#2] at (\x,\y) {};}
  \newcommand{\groups}{\lines{\bondsKept}{bond} \nodes{\sitesAl}{al} \nodes{\keptMg}{mg} \nodes{\keptZn}{zn} \nodes{\keptX}{va}}
  \newcommand{\dropped}[1]{%
    \begin{scope}[transparency group, opacity=#1]
      \lines{\bondsDropped}{bond} \nodes{\droppedMg}{mg} \nodes{\droppedZn}{zn}
    \end{scope}}
  % ---------------- (a) connected ----------------
  \begin{scope}
    \groups \dropped{1}
    \node[above lab] at (-0.45,0.62) {\panel{a}};
  \end{scope}
  % ---------------- (b) dropped ----------------
  \begin{scope}[shift={(9.9,0)}]
    \lines{\bondsKept}{rim} \lines{\bondsKept}{region}
    \groups \dropped{0.3}
    \node[above lab] at (-0.45,0.62) {\panel{b}};
  \end{scope}
  % ---------------- (c) solvent atoms ----------------
  \begin{scope}[shift={(19.8,0)}]
    \lines{\bondsClusters}{rim} \lines{\bondsClusters}{region}
    \groups \dropped{0.3}
    \nodes{\joinedAl}{joined}
    \foreach \x/\y/\n in \counts \node[number] at (\x,\y) {\n};
    \draw[pointer] (2.5,0.62) -- (2.5,-0.35);
    \node[anchor=south, inner sep=1.5pt, font=\footnotesize] at (2.9,0.62) {\texttt{cluster\char95 id} 1};
    \draw[pointer] (6,-6.70) -- (6,-5.71);
    \node[anchor=north, inner sep=1.5pt, font=\footnotesize] at (6,-6.70) {\texttt{cluster\char95 id} 2};
    \node[above lab] at (-0.45,0.62) {\panel{c}};
  \end{scope}
  % ---------------- legend ----------------
  \begin{scope}[shift={(-0.4,-8.2)}, x=1cm, y=1cm]
    \node[al] at (0.25,0) {}; \node[leg] at (0.5,0) {Al (solvent)};
    \node[mg] at (2.95,0) {}; \node[leg] at (3.2,0) {Mg};
    \node[zn] at (4.25,0) {}; \node[leg] at (4.5,0) {Zn};
    \node[va] at (5.5,0) {}; \node[leg] at (5.75,0) {vacancy};
    \draw[bond] (7.45,0) -- ++(0.55,0); \node[leg] at (8.1,0) {nearest neighbors};
    \begin{scope}[transparency group, opacity=0.3] \node[mg] at (11.4,0) {}; \end{scope} \node[leg] at (11.65,0) {dropped};
    \node[al] at (13.45,0) {}; \node[joined] at (13.45,0) {}; \node[leg] at (13.7,0) {joins the group};
  \end{scope}
\end{tikzpicture}
\end{document}
